\documentclass[11pt]{article}
\usepackage[a4paper,margin=25mm]{geometry}
\usepackage{amsmath,amssymb,amsthm}
\usepackage[hidelinks]{hyperref}
\emergencystretch=2em
\newtheorem{theorem}{Theorem}
\newcommand{\F}{\mathbb F}
\title{One-variable Veronese rings have F-signature one\newline
\large Disproof of Conjecture 00000003904}
\author{gaochengzhecpu}\date{}
\begin{document}\maketitle
\begin{abstract}
For every $d>0$, the Veronese subring $V_d=\F_2[X^d]$
is a polynomial ring in one variable. We explicitly compute its
Frobenius pushforward as a free module of rank $2^e$ at level $e$.
Its actual maximal splitting counts, dimension-one normalization
and F-signature therefore give $s(V_d)=1$ for every $d>0$.
The proposed $d^{-1}(1+c/d+O(d^{-2}))$ tends to zero, so it cannot
hold. The accompanying Lean project proves the module computation,
Krull dimension and asymptotic contradiction, rather than assuming
a general regular-ring F-signature formula.
\end{abstract}

\section*{The permitted one-variable instance}
The source defines the $d$-Veronese subring by all degree-$d$
monomials and asserts
\[
 s(V_d)=d^{-n}\bigl(1+c(n)/d+O(d^{-2})\bigr).
\]
Neither source language excludes $n=1$. Take the actual perfect
field $\F_2$ and a single variable $X$. The degree-$d$ monomials
generate exactly $V_d=\F_2[X^d]\subseteq\F_2[X]$.
The map
\[
 \phi_d:\F_2[T]\longrightarrow V_d,\qquad f(T)\longmapsto f(X^d)
\]
is onto by the definition of the subalgebra and injective for $d>0$:
distinct exponents $j$ give distinct exponents $dj$.
Thus it is an algebra isomorphism. In particular $V_d$ is
Noetherian and has Krull dimension one. The formal dimension
proof uses the polynomial PID property for the upper bound and
the field-to-polynomial dimension lower bound for the reverse inequality.

We use the standard graded F-signature definition~\cite{Singh},
with the homogeneous maximal ideal understood. If $a_e$ is the
maximal rank of a free direct summand of $F_*^e V_d$, then
\[
 s(V_d)=\lim_{e\to\infty}\frac{a_e}{q^{\dim V_d}},\qquad q=2^e.
\]
The coefficient field is perfect, so no residue-field correction is
needed. We compute the actual entire module below.

\newpage
\section*{A complete basis for the actual Frobenius module}
As an additive group, $F_*^e V_d$ is $V_d$ itself; its scalar
action is twisted:
\[
 a\cdot m=a^q m,\qquad q=2^e.
\]
This is restriction of scalars along the actual iterated Frobenius,
not the ordinary multiplication action.
Over $\F_2$, coefficients satisfy $a^q=a$, hence
$f(T)^q=f(T^q)$ for every polynomial $f$.

For the polynomial ring $B=\F_2[T]$, define
\[
 S_q:B^q\longrightarrow F_*^e B,\qquad
 (f_0,\ldots,f_{q-1})\longmapsto
 \sum_{i=0}^{q-1} f_i(T^q)T^i.
\]
It is $B$-linear for the twisted action. Its coefficient at
$qn+i$, with $0\leq i<q$, is exactly the coefficient of $T^n$
in $f_i$. Distinct residues cannot contribute to this coefficient,
so $S_q$ is injective. Conversely, the monomial $aT^N$ is the
image of the tuple having $aT^{\lfloor N/q\rfloor}$ in position
$N\bmod q$ and zeros elsewhere. The quotient-remainder identity
$N=q\lfloor N/q\rfloor+(N\bmod q)$ proves surjectivity for every
monomial, and addition proves it for all polynomials.

Transporting this linear equivalence through $\phi_d$ gives the
actual equivalence
\[
 V_d^{\,q}\ \cong_{V_d}\ F_*^e V_d.
\]
The standard basis vectors are sent to the actual polynomials
\[
 1,X^d,X^{2d},\ldots,X^{(q-1)d}.
\]
This includes $e=0$, where $q=1$, and makes no coprimality
assumption between $d$ and the characteristic.

\section*{Maximal free summands, not just a lower bound}
Define the free-summand ranks of an $A$-module $M$ to be the
natural numbers $k$ for which there are actual $A$-linear maps
\[
 i:A^k\longrightarrow M,\qquad r:M\longrightarrow A^k,
 \qquad r\circ i=\operatorname{id}.
\]
Their supremum is the free rank. The proved equivalence above
exhibits the free rank $q$. For the upper bound, every such split
injection is injective, and the strong rank condition for the
commutative domain $V_d$ gives $k\leq q$. Therefore the actual
maximal free-summand rank is $a_e=q$, not merely at least $q$.

\newpage
\section*{The limit and the asymptotic contradiction}
The proved dimension and splitting computation give
\[
 \frac{a_e}{q^{\dim V_d}}=\frac{2^e}{2^e}=1
 \qquad(d>0,\ e\geq0).
\]
Consequently the genuine splitting-count sequence has limit one.

\begin{theorem}
For every $d>0$, the one-variable Veronese ring
$V_d=\F_2[X^d]$ has graded F-signature $s(V_d)=1$.
There are no real $c$ and remainder $r(d)=O(d^{-2})$ such that,
for all sufficiently large positive integers $d$,
\[
 s(V_d)=d^{-1}\bigl(1+c/d+r(d)\bigr).
\]
\end{theorem}
\begin{proof}
The first assertion is the module and dimension computation above.
For the second, a genuine $O(d^{-2})$ remainder tends to zero.
Thus the displayed right side tends to zero, whereas the left side
is identically one. Uniqueness of a real limit gives a contradiction.
\end{proof}

This refutes the source's F-signature asymptotic, hence the
asserted conjunction. It makes no claim that the strong
F-regularity clause fails. A single permitted characteristic and
variable count suffice to refute the universal formula.

\section*{Formalization scope}
The Lean project uses the actual polynomial subalgebra generated by
all degree-$d$ monomials and proves its equality to the singleton
adjoin. Injective polynomial expansion gives the algebra isomorphism.
The Frobenius module is built with restriction of scalars along
Mathlib's \texttt{iterateFrobenius}; its displayed scalar action is
checked explicitly. Polynomial coefficients prove the complete
residue decomposition, and a further theorem identifies every
transported basis vector with $X^{di}$.

The free rank is defined by actual split linear maps, with boundedness
and attainment proved. Mathlib's actual Krull dimension is proved
one. F-signature is defined as the limit of these actual normalized
counts, and the limit is proved to exist and equal one. The final
theorem quantifies over every real correction constant and every
remainder satisfying Mathlib's actual \texttt{IsBigO}; it does not
replace asymptotic notation by a pointwise error condition.

The standard graded definition is the only terminology bridge to
the source; see~\cite{Singh}, p.~1. The computation does not use
an external F-signature evaluation theorem. No Lean theorem
comparing this graded definition with localization is claimed.

\newpage
\section*{Reproduction and review}
The complete project pins Lean 4.19.0. The Mathlib revision is
\begin{center}
\texttt{c44e0c8ee63ca166450922a373c7409c5d26b00b}.
\end{center}
From the \texttt{lean/} directory run:
\begin{verbatim}
lake exe cache get
lake build
lake env lean -DwarningAsError=true Main.lean
\end{verbatim}
Compile the paper with \texttt{tectonic main.tex}.
The source includes explicit axiom reports for the ring,
Frobenius basis, dimension, free-summand computation and final
asymptotic theorem. The only printed dependencies are standard
foundational axioms. Actual logs, source hashes and fresh-build
evidence accompany the submission.

All theorems computing the ring, dimension and signature require
$d>0$. Definitions at $d=0$ are merely total Lean definitions;
the asymptotic contradiction uses only eventually positive $d$.
No finite computation is an oracle for the formal proof.

The work received author, parent and internal cross-review.
No external independent review is claimed. The complete proof
rests on an explicit free-module decomposition and a limit, so its
conclusion does not depend on trusting a numerical estimate.

\begin{thebibliography}{9}
\bibitem{source} The Last Math Competition,
Conjecture 00000003904. Exact bilingual source included.
\bibitem{Singh} A. K. Singh,
\emph{The F-signature of an affine semigroup ring} (2004),
\href{https://arxiv.org/abs/math/0409143}{arXiv:math/0409143},
p.~1, definition of F-signature for local or graded rings.
\end{thebibliography}
\end{document}
